\daytitle{6}{牛顿运动定律}

\begin{learningbox}{本节导读}
牛顿三大定律是整个经典力学的基石。本节把受力分析和运动学接起来：先把研究对象受哪些真实力画清楚，再把合力写成 $F_{\text{合}} = ma$。读本节时要始终检查三件事：研究对象是谁、合力沿哪个方向、加速度与运动学量怎样衔接。
\end{learningbox}

\textbf{本节读题路线：} 先选研究对象和坐标轴，再画真实受力图；沿每个方向列 $F_{\text{合}}=ma$，求出加速度或未知力后，再接运动学公式。作用力与反作用力作用在不同物体上，不要画进同一个受力图里。

\section*{一、学习目标}

\begin{enumerate}
  \item 能用惯性定律说明：力不是维持运动的原因，而是改变运动状态的原因。
  \item 能在受力图中找出合力，并用 $F_{\text{合}} = ma$ 建立方程。
  \item 能区分作用力与反作用力、平衡力，并说明它们的对象不同。
  \item 会用正交分解法列动力学方程。
  \item 能在水平面、斜面、连接体和竖直升降问题中，先画受力图，再按方向列出合力与加速度方程。
\end{enumerate}

\section*{二、初高中衔接}

初中对牛顿定律的学习停留在定性层面：牛顿第一定律叫“惯性定律”，牛顿第三定律叫“力的作用是相互的”。

高中在以下方面深化：
\begin{enumerate}
  \item \textbf{牛顿第二定律定量化：} $F_{\text{合}} = ma$，已知力可以算加速度，已知加速度可以反推力。
  \item \textbf{力学与运动学打通：} 牛顿定律是“桥梁”——受力分析 $\to$ 合力 $\to$ 加速度 $\to$ 运动学公式。
  \item \textbf{多物体系统：} 整体法和隔离法分析连接体问题。
\end{enumerate}

\textbf{科学发展视角：} 牛顿第二定律是许多控制系统的核心。电梯启动、汽车制动、无人机悬停、机械臂加速，都要把“受力”转化为“加速度”，再预测下一刻的位置和速度。现代设备常由传感器测量运动状态，再由控制程序反推需要施加多大的力。

\section*{三、核心概念}

\begin{conceptbox}{牛顿第一定律（惯性定律）}

\textbf{内容：} 在惯性参考系中，一切物体在不受外力作用（或合力为零）时，总保持匀速直线运动状态或静止状态。

\textbf{惯性：} 物体保持原有运动状态不变的性质。
\begin{itemize}
  \item 惯性是物体的固有属性，与是否受力、是否运动无关。
  \item 质量是惯性大小的量度——质量越大，惯性越大，运动状态越难改变。
\end{itemize}

\textbf{伽利略的理想实验：} 斜面实验表明，如果没有摩擦力，小球将沿另一斜面上升到相同高度。若另一斜面倾角为零（水平），小球将一直运动下去——推翻了亚里士多德“力是维持运动的原因”的错误观点。

\begin{center}
% TikZ 图：伽利略理想斜面实验
% Day 06：牛顿运动定律
\begin{tikzpicture}[scale=0.82,>=Latex]
  % 左图：右侧仍为斜面，小球能上升到与初始位置相同的高度。
  \begin{scope}[shift={(-3.75,0)}]
    \node[font=\footnotesize] at (0,1.95) {(1) 右侧仍为斜面};
    \draw[thick] (-2.8,1.35) -- (0,0) -- (2.8,1.35);
    \draw[dashed,gray] (-2.95,1.35) -- (2.95,1.35);
    \draw[<->,gray] (-3.15,0) -- (-3.15,1.35) node[midway,left] {$h$};

    \filldraw[fill=white,draw=black,thick] (-2.25,1.08) circle (0.13);
    \filldraw[fill=white,draw=black,thick] (2.25,1.08) circle (0.13);
    \draw[->,thick,red] (-1.95,0.93) -- (-1.18,0.56) node[midway,below] {$v$};
    \draw[->,thick,red] (1.15,0.56) -- (1.92,0.93);

    \node[font=\footnotesize] at (0,-0.42) {无摩擦：上升到同一高度};
  \end{scope}

  % 右图：右侧斜面逐渐放平，极限情形为水平面上的匀速直线运动。
  \begin{scope}[shift={(3.75,0)}]
    \node[font=\footnotesize] at (0,1.95) {(2) 右侧放平成水平面};
    \draw[thick] (-2.8,1.35) -- (0,0) -- (3.25,0);
    \draw[dashed,gray] (-2.95,1.35) -- (-1.35,1.35);
    \draw[<->,gray] (-3.15,0) -- (-3.15,1.35) node[midway,left] {$h$};

    \filldraw[fill=white,draw=black,thick] (-2.25,1.08) circle (0.13);
    \filldraw[fill=white,draw=black,thick] (0.75,0.13) circle (0.13);
    \draw[->,thick,red] (-1.95,0.93) -- (-1.18,0.56) node[midway,below] {$v$};
    \draw[->,thick,red] (0.98,0.13) -- (2.28,0.13) node[midway,above] {$v$};

    \node[font=\footnotesize] at (0,-0.42) {无摩擦：速度大小保持};
  \end{scope}
\end{tikzpicture}


\figcaption{伽利略理想斜面：无摩擦时，运动不需要持续外力维持}\label{fig:newton-first-ideal}
\end{center}

\textbf{科学方法提醒：} 牛顿第一定律不能直接用实验验证（无法完全消除外力），它是由实验和推理得出的。它同时揭示了什么样的参考系可近似看作惯性参考系：在这样的参考系中，不受合外力的物体会保持静止或匀速直线运动。

\end{conceptbox}

\begin{conceptbox}{牛顿第二定律}

\textbf{内容：} 物体的加速度与所受合力成正比，与质量成反比，加速度方向与合力方向相同。

\[
\boxed{F_{\text{合}} = ma}
\]

\begin{itemize}
  \item \textbf{矢量性：} $F_{\text{合}}$ 和 $a$ 方向相同。
  \item \textbf{瞬时性：} 在高中牛顿力学模型中，合力一旦改变，加速度立即随之改变（同时性）。
  \item \textbf{叠加性：} 各个力的作用可以按矢量叠加，最终由合力决定物体的加速度。
\end{itemize}

\textbf{解题关键：} 用牛顿第二定律建立方程，把受力分析的结果代入。
\[
\text{受力分析} \;\to\; \text{求合力}\; F_{\text{合}} \;\to\; F_{\text{合}} = ma \;\to\; \text{求加速度或某个力}
\]

\textbf{方法流程：} 牛顿第二定律的核心是把受力图转化为方程。每一步都要保留研究对象和正方向，最后再写出合力与加速度的关系。

\begin{center}
\begin{tikzpicture}[>=stealth,node distance=1.35cm]
  \tikzstyle{stepbox}=[draw,thick,minimum width=2.35cm,minimum height=0.82cm,align=center]

  \node[stepbox] (obj) {选对象};
  \node[stepbox,right=of obj] (force) {画真实力};
  \node[stepbox,right=of force] (axis) {定正方向};
  \node[stepbox,right=of axis] (eq) {列方程};

  \draw[->,thick] (obj) -- (force);
  \draw[->,thick] (force) -- (axis);
  \draw[->,thick] (axis) -- (eq);

  \node[below=0.48cm of obj,align=center] {隔离研究对象};
  \node[below=0.48cm of force,align=center] {重力、弹力\\摩擦力};
  \node[below=0.48cm of axis,align=center] {把力分解\\带正负号};
  \node[below=0.48cm of eq,align=center] {$\sum F_x=ma_x$\\$\sum F_y=ma_y$};
\end{tikzpicture}


\figcaption{动力学题的标准流程：从研究对象到方程}\label{fig:newton-solution-flow}
\end{center}

\begin{center}
% TikZ 图：牛顿第二定律示意（水平面上）
% Day 06：牛顿运动定律
\begin{tikzpicture}[scale=0.9,>=Latex]
  \draw[thick] (-1,0) -- (5,0);
  % 物体
  \draw[thick] (0.5,0) rectangle (2.5,1.5);
  \node at (1.5,0.75) {$m$};
  % 拉力
  \draw[->,thick,red] (2.5,0.75) -- (4.5,0.75) node[above] {$F$};
  % 摩擦力
  \draw[->,thick,red] (0.5,0.75) -- (-1,0.75) node[above] {$f$};
  % 加速度箭头在上方
  \draw[->,ultra thick,blue] (1.5,2.5) -- (3.5,2.5) node[above] {$a = \dfrac{F-f}{m}$};
\end{tikzpicture}


\figcaption{水平面滑动时 $F_{\text{合}} = F - f_k = ma$}\label{fig:newton-second}
\end{center}

\end{conceptbox}

\begin{conceptbox}{牛顿第三定律}

\textbf{内容：} 两个物体之间的作用力和反作用力总是大小相等、方向相反、作用在同一条直线上。

\[
\vect{F}_{AB} = -\vect{F}_{BA}
\]

\textbf{四个特征：}
\begin{enumerate}
  \item \textbf{等大：} $|F_{AB}| = |F_{BA}|$。
  \item \textbf{反向：} 方向相反。
  \item \textbf{共线：} 作用在同一直线上。
  \item \textbf{同时：} 同生同灭，没有先后。
\end{enumerate}

\textbf{关键区别：} 作用力与反作用力（牛顿第三定律）$\neq$ 平衡力（牛顿第一定律）。

\begin{center}
\begin{tabular}{lll}
\toprule
 & 作用力与反作用力 & 平衡力 \\
\midrule
对象 & 两个不同物体 & 同一个物体 \\
性质 & 一定相同 & 可以不同 \\
时间 & 同生同灭 & 无此要求 \\
效果 & 各自产生效果 & 合力为零 \\
\bottomrule
\end{tabular}
\end{center}

\begin{center}
% TikZ 图：作用力与反作用力、平衡力对比
% Day 06：牛顿运动定律
\begin{tikzpicture}[scale=0.9,>=Latex]
  % 左图：作用力与反作用力作用在不同物体上。
  \begin{scope}
    \node[font=\bfseries] at (1.45,1.75) {作用力与反作用力};
    \draw[thick] (0,0) rectangle (1.1,0.75);
    \draw[thick] (1.8,0) rectangle (2.9,0.75);
    \node at (0.55,0.38) {$A$};
    \node at (2.35,0.38) {$B$};
    \draw[->,thick,blue] (1.1,0.38) -- (0.12,0.38) node[midway,above] {\footnotesize $B$ 对 $A$};
    \draw[->,thick,red] (1.8,0.38) -- (2.78,0.38) node[midway,above] {\footnotesize $A$ 对 $B$};
    \node[font=\footnotesize,align=center] at (1.45,-0.55) {等大反向、共线\\但作用在两个物体上};
  \end{scope}

  % 右图：平衡力作用在同一个物体上。
  \begin{scope}[shift={(5.2,0)}]
    \node[font=\bfseries] at (1.05,1.75) {平衡力};
    \draw[thick] (-0.2,-0.05) -- (2.3,-0.05);
    \draw[thick,fill=white] (0.45,0) rectangle (1.65,0.75);
    \node at (1.05,0.38) {$m$};
    \draw[->,thick,blue] (1.05,0.75) -- (1.05,1.45) node[right] {$N$};
    \draw[->,thick,red] (1.05,0) -- (1.05,-0.82) node[right] {$G$};
    \node[font=\footnotesize,align=center] at (1.05,-1.18) {作用在同一物体上\\合力可以为零};
  \end{scope}
\end{tikzpicture}


\figcaption{作用力与反作用力和平衡力的研究对象不同}\label{fig:newton-third-compare}
\end{center}

\end{conceptbox}

\section*{四、公式与适用条件}

\textbf{列式前检查：}
\begin{enumerate}
  \item 研究对象是单个物体还是整体？不同对象不能共用同一张受力图。
  \item 是否选了惯性参考系？若选加速参考系，必须另行说明惯性力。
  \item 坐标轴是否沿运动方向、斜面方向或其他方便方向选好？
  \item 每个方向的加速度是否真的为零，不要把“没有运动”误判成“没有受力”。
\end{enumerate}

\begin{learningbox}{从受力图到方程的三行分析清单}
牛顿第二定律题的关键不是把公式背熟，而是把受力图翻译成方程。每选一个研究对象，就按下面三行写：

\begin{center}
\begin{tabular}{lll}
\toprule
步骤 & 要写清什么 & 例子 \\
\midrule
1 & 选轴与正方向 & 沿斜面向下为 $x$ 轴正方向 \\
2 & 写该方向所有力的代数和 & $mg\sin\theta-f_k$ \\
3 & 写该方向的加速度 & $a_x=a$，垂直斜面 $a_y=0$ \\
\bottomrule
\end{tabular}
\end{center}

于是沿斜面方向写成
\[
mg\sin\theta-f_k=ma,
\]
垂直斜面方向写成
\[
N-mg\cos\theta=0.
\]
\textbf{判断标准：} 方程左边必须来自受力图中的真实力或真实力的分量；右边是同一方向的 $ma$。若某方向没有加速度，右边才写 $0$，不是因为“这个方向没有运动”就自动没有受力。
\end{learningbox}

\subsection*{核心公式}

\[
\boxed{F_{\text{合}} = ma}
\]

\begin{itemize}
  \item 适用条件：宏观、低速（远小于光速）、惯性参考系。
  \item $F_{\text{合}}$ 是物体所受所有力的矢量和，不是某个力。
  \item $a$ 与 $F_{\text{合}}$ 瞬时对应。
\end{itemize}

\textbf{参考系提醒：} 本讲义默认在地面等近似惯性参考系中列牛顿第二定律。若题目主动要求“站在加速车厢中看”，就进入非惯性参考系，需要额外处理惯性力；初学阶段优先选惯性参考系，避免把惯性力当成真实相互作用力。

\subsection*{正交分解法}

当物体受力不在同一直线上时，把力分解到两个互相垂直的方向上。分解前先规定正方向；列方程时把每个力的分量带正负号代入。

\begin{center}
% TikZ 图：正交分解
% Day 06：牛顿运动定律
\begin{tikzpicture}[scale=1,>=Latex]
  \draw[->,thick] (0,0) -- (3,0) node[right] {$x$};
  \draw[->,thick] (0,0) -- (0,3) node[above] {$y$};
  % 力 F（斜向上）
  \draw[->,ultra thick,red] (0,0) -- (2.5,1.5) node[above right] {$F$};
  % 分力
  \draw[dashed,thick,red] (0,0) -- (2.5,0);
  \draw[dashed,thick,red] (2.5,0) -- (2.5,1.5);
  % 角度
  \draw (0.8,0) arc (0:31:0.8);
  \node at (1.1,0.2) {$\theta$};
  % 标注分量
  \node[red,below] at (1.25,0) {$F_x = F\cos\theta$};
  \node[red,right] at (2.5,0.75) {$F_y = F\sin\theta$};
\end{tikzpicture}


\figcaption{正交分解：$F_x = F\cos\theta$，$F_y = F\sin\theta$}\label{fig:orthogonal-decomp}
\end{center}

\noindent 列方程标准形式：
\[
\begin{cases}
F_{x,\text{合}} = ma_x \\
F_{y,\text{合}} = ma_y
\end{cases}
\]

若已知物体被约束在某一方向运动，常把该方向取为 $x$ 轴；垂直于运动方向没有加速度时，才有 $a_y = 0$，从而 $y$ 方向合力为零。若物体做曲线运动或竖直方向也在加速，就不能直接写 $a_y=0$。

\section*{五、典型模型}

\subsection*{模型 1：水平面上的动力学}

物体在水平拉力 $F$ 作用下已经在水平面上滑动：
\[
F - f_k = ma,\qquad N = mg,\qquad f_k=\mu_k N
\]
若拉力 $F$ 斜向上、与水平方向成 $\theta$ 角，且物体没有离开地面，则 $N = mg - F\sin\theta$，滑动摩擦力 $f_k=\mu_k N$：
\[
F\cos\theta - \mu_k (mg - F\sin\theta) = ma
\]

\subsection*{模型 2：斜面上的动力学}

物体沿斜面\textbf{向下滑}（沿斜面向下为正方向）：
\[
mg\sin\theta - f_k = ma,\qquad N = mg\cos\theta,\qquad f_k = \mu_k N = \mu_k mg\cos\theta
\]
\[
a_{\text{下滑}} = g\sin\theta - \mu_k g\cos\theta
\]

物体沿斜面\textbf{向上滑}（沿斜面向上为正方向）：
\[
-mg\sin\theta - f_k = ma,\qquad N = mg\cos\theta,\qquad f_k = \mu_k N
\]
\[
a_{\text{上滑}} = -\,(g\sin\theta + \mu_k g\cos\theta)
\]
（负号表示加速度沿斜面向下，大小为 $g\sin\theta + \mu_k g\cos\theta$）

画斜面受力图时，先画真实力：重力、支持力、摩擦力。重力的两个分力只是列方程时的计算工具，不是额外多出来的力。

\begin{center}
% TikZ 图：斜面上牛顿第二定律（下滑时）
% Day 06：牛顿运动定律
% 左高右低斜面，受真实力（G/N/f）+ 重力分力（虚线）
\begin{tikzpicture}[scale=0.9,>=Latex]
  % 斜面
  \draw[thick] (0,2.5) -- (5,0.5);
  \draw[thick] (0,0.5) -- (5,0.5);
  % 角度
  \draw (0.8,0.5) arc (0:-21.8:0.8);
  \node at (1.1,0.35) {$\theta$};
  
  % 沿斜面画方块
  \begin{scope}[shift={(2.5,1.5)},rotate=-21.8]
    \draw[thick] (-0.6,-0.5) rectangle (0.6,0.5);
    \node at (0,0) {$m$};
    
    % 真实力
    \draw[->,thick,blue] (0,0.5) -- (0,1.4) node[above] {$N$};
    \draw[->,thick,red] (-0.6,0) -- (-1.4,0) node[above] {$f$};
    % 加速度 a（沿 x 正方向，即沿斜面向下）
    \draw[->,ultra thick,red] (0.6,0) -- (1.6,0) node[above] {$a$};
    
    % 重力分力（虚线，橙色）
    \draw[->,thick,dashed,orange] (0,0) -- (1.4,0) node[right] {$G\sin\theta$};
    \draw[->,thick,dashed,orange] (0,0) -- (0,-0.9) node[below] {$G\cos\theta$};
  \end{scope}
  
  % G（全局坐标系）
  \draw[->,thick,blue] (2.5,1.0) -- (2.5,-0.1) node[below] {$G$};
  
  % 条件
  \node[below,font=\footnotesize] at (2.5,-0.4) {（下滑时 $f$ 沿斜面向上，虚线为重力分力）};
\end{tikzpicture}


\nopagebreak
\figcaption{斜面上物体的受力与分解：实线为真实力，虚线为重力分力}\label{fig:newton-incline}
\end{center}

\subsection*{模型 3：连接体问题（基础）}

两个物体用轻绳连接，在水平拉力作用下一起沿水平面运动。若水平方向的外力只有拉力和各物体受到的摩擦力，先用整体法求加速度：
\[
F - f_{\text{总}} = (m_1 + m_2) a
\]
再用隔离法求绳中张力。

\begin{center}
% TikZ 图：连接体整体法与隔离法
% Day 06：牛顿运动定律
\begin{tikzpicture}[scale=0.9,>=Latex]
  \draw[thick] (-0.8,0) -- (6.0,0);

  % 两个物体与轻绳：外力拉右侧 A，左侧 B 被绳拉动。
  \draw[thick,fill=white] (0.2,0) rectangle (1.6,0.9);
  \draw[thick,fill=white] (2.5,0) rectangle (3.9,0.9);
  \draw[thick] (1.6,0.45) -- (2.5,0.45);
  \node at (0.9,0.45) {$B$};
  \node at (3.2,0.45) {$A$};
  \node[font=\footnotesize] at (2.05,0.7) {轻绳};

  % 外力、加速度与隔离 B 时的张力和摩擦力。
  \draw[->,ultra thick,red] (3.9,0.45) -- (5.05,0.45) node[midway,above] {$F$};
  \draw[->,thick,blue] (1.6,0.45) -- (2.38,0.45) node[midway,below] {$T$};
  \draw[->,thick,red] (0.2,0.28) -- (-0.45,0.28) node[midway,above] {$f_B$};
  \draw[->,thick,blue] (1.55,1.35) -- (3.0,1.35) node[midway,above] {$a$};

  % 方法提示。
  \draw[dashed,rounded corners=2pt] (-0.05,-0.15) rectangle (4.15,1.1);
  \node[font=\footnotesize,align=center] at (2.05,-0.62)
    {整体法：只算外力求 $a$；隔离 $B$：用 $T-f_B=m_Ba$ 求张力};
\end{tikzpicture}


\figcaption{连接体：整体法求共同加速度，隔离法求内力}\label{fig:connected-body-newton}
\end{center}

\subsection*{模型 4：超重与失重}

超重和失重说的是\textbf{视重}变化，不是重力 $mg$ 本身改变。以升降机中的人为例，支持力 $N$ 是人对秤的压力读数对应的力，也是人的视重。

\begin{center}
% TikZ 图：升降机中的超重、正常视重与失重
% Day 06：牛顿运动定律
\begin{tikzpicture}[scale=0.84,>=Latex,line cap=round,line join=round]
  % 向上加速：N 的箭头明显长于 mg。
  \begin{scope}[xshift=0cm]
    \node[font=\small] at (0,2.25) {向上加速};
    \draw[thick] (-1.0,-1.20) rectangle (1.0,1.20);
    \draw[thick] (-0.45,-0.75) rectangle (0.45,-0.10);
    \node at (0,-0.42) {$m$};
    \draw[->,very thick,blue] (0,-0.10) -- (0,1.02)
      node[right] {$N$};
    \draw[->,very thick,red] (0,-0.75) -- (0,-1.52)
      node[right] {$mg$};
    \draw[->,ultra thick] (-0.74,1.35) -- (-0.74,2.00)
      node[left] {$a$};
    \node[font=\small] at (0,-2.05) {$N>mg$};
  \end{scope}

  % 匀速或静止：两箭头等长。
  \begin{scope}[xshift=4cm]
    \node[font=\small] at (0,2.25) {匀速或静止};
    \draw[thick] (-1.0,-1.20) rectangle (1.0,1.20);
    \draw[thick] (-0.45,-0.75) rectangle (0.45,-0.10);
    \node at (0,-0.42) {$m$};
    \draw[->,very thick,blue] (0,-0.10) -- (0,0.78)
      node[right] {$N$};
    \draw[->,very thick,red] (0,-0.75) -- (0,-1.63)
      node[right] {$mg$};
    \node[font=\small] at (0,-2.05) {$N=mg$};
  \end{scope}

  % 向下加速：N 的箭头明显短于 mg。
  \begin{scope}[xshift=8cm]
    \node[font=\small] at (0,2.25) {向下加速};
    \draw[thick] (-1.0,-1.20) rectangle (1.0,1.20);
    \draw[thick] (-0.45,-0.75) rectangle (0.45,-0.10);
    \node at (0,-0.42) {$m$};
    \draw[->,very thick,blue] (0,-0.10) -- (0,0.52)
      node[right] {$N$};
    \draw[->,very thick,red] (0,-0.75) -- (0,-1.82)
      node[right] {$mg$};
    \draw[->,ultra thick] (-0.74,2.00) -- (-0.74,1.35)
      node[left] {$a$};
    \node[font=\small] at (0,-2.05) {$N<mg$};
  \end{scope}
\end{tikzpicture}


\figcaption{升降机中的视重：支持力 $N$ 随竖直加速度改变}\label{fig:apparent-weight-elevator}
\end{center}

取竖直向上为正方向：
\[
N-mg=ma,\qquad N=m(g+a).
\]

\begin{itemize}
  \item 向上加速或向下减速：$a>0$，$N>mg$，表现为超重。
  \item 向下加速或向上减速：$a<0$，$N<mg$，表现为失重。
  \item 自由下落：$a=-g$，$N=0$，表现为完全失重。
\end{itemize}

\newpage
\section*{六、例题精讲}

\begin{examplebox}{水平面上的牛顿第二定律}

\textbf{题意分析：} 一个 $m = 4\ \mathrm{kg}$ 的物体在水平面上，受水平拉力 $F = 20\ \mathrm{N}$，物体与地面间的动摩擦因数 $\mu_k = 0.3$，$g = 10\ \mathrm{m/s^2}$。求物体的加速度。

\textbf{物理模型：} 水平面动力学，受拉力 $F$ 和滑动摩擦力 $f_k$。

\textbf{解题过程：}

正压力 $N = mg = 4 \times 10 = 40\ \mathrm{N}$
滑动摩擦力 $f_k = \mu_k N = 0.3 \times 40 = 12\ \mathrm{N}$

水平方向合力：$F_{\text{合}} = F - f_k = 20 - 12 = 8\ \mathrm{N}$

由牛顿第二定律：
\[
a = \frac{F_{\text{合}}}{m} = \frac{8}{4} = 2\ \mathrm{m/s^2}
\]

\textbf{方法总结：} 本题拉力水平，竖直方向没有其他分量，所以 $N = mg$；再算 $f_k = \mu_k N$，最后用 $F - f_k = ma$。若拉力或推力有竖直分量，必须重新求 $N$。

\end{examplebox}

\begin{examplebox}{斜面上的牛顿第二定律}

\textbf{题意分析：} 质量为 $m = 2\ \mathrm{kg}$ 的物体从倾角 $\theta = 37^\circ$（$\sin37^\circ = 0.6$，$\cos37^\circ = 0.8$）的斜面顶端由静止下滑，物体与斜面间的动摩擦因数 $\mu_k = 0.25$，$g = 10\ \mathrm{m/s^2}$。求：
\begin{enumerate}
  \item 物体下滑的加速度。
  \item 若斜面长 $L = 5\ \mathrm{m}$，物体滑到底端所需时间和末速度。
\end{enumerate}

\textbf{物理模型：} 斜面动力学 $\to$ 匀加速直线运动。

\textbf{解题过程：}

\begin{enumerate}
  \item 沿斜面方向合力：
  \[
  mg\sin\theta - \mu_k mg\cos\theta = ma
  \]
  \[
  a = g(\sin\theta - \mu_k\cos\theta) = 10 \times (0.6 - 0.25 \times 0.8) = 10 \times 0.4 = 4\ \mathrm{m/s^2}
  \]
  
  \item 由运动学公式（$v_0 = 0$）：
  \[
  L = \frac12 a t^2,\qquad 5 = \frac12 \times 4 \times t^2,\qquad t^2 = 2.5,\qquad t \approx 1.58\ \mathrm{s}
  \]
  \[
  v = a t = 4 \times 1.58 \approx 6.32\ \mathrm{m/s}
  \]
\end{enumerate}

\textbf{方法总结：} 斜面动力学先沿斜面、垂直斜面建轴，再把重力分解为 $mg\sin\theta$ 和 $mg\cos\theta$。沿斜面分量决定加速度，垂直斜面分量决定支持力。

\end{examplebox}

\begin{examplebox}{连接体问题}

\textbf{题意分析：} 质量 $m_A = 3\ \mathrm{kg}$、$m_B = 2\ \mathrm{kg}$ 的两个物体用轻绳连接，放在光滑水平面上。$B$ 在左、$A$ 在右，在 $A$ 上施加水平向右的拉力 $F = 10\ \mathrm{N}$。求：
\begin{enumerate}
  \item 整体的加速度。
  \item 绳中张力 $T$。
\end{enumerate}

\textbf{物理模型：} 连接体，光滑无摩擦。

\textbf{解题过程：}

\begin{enumerate}
  \item 整体法（把 $A$、$B$ 看作一个整体）：
  \[
  F = (m_A + m_B) a,\qquad 10 = 5a,\qquad a = 2\ \mathrm{m/s^2}
  \]
  
  \item 隔离 $B$（只分析 $B$ 受的力）：
  
  $B$ 只受绳的拉力 $T$，由牛顿第二定律：
  \[
  T = m_B a = 2 \times 2 = 4\ \mathrm{N}
  \]
  
  验证（隔离 $A$）：$F - T = m_A a$，$10 - 4 = 6 = 3 \times 2$，正确。
\end{enumerate}

\textbf{方法总结：}
\begin{enumerate}
  \item 连接体问题先用整体法求加速度，再用隔离法求内力（张力）。
  \item 整体法只考虑整体受的外力，不考虑内部相互作用力。
  \item 光滑表面摩擦可忽略，粗糙表面需要加上整体摩擦力。
\end{enumerate}

\end{examplebox}

\begin{examplebox}{升降机中的超重与失重}

\textbf{题意分析：} 质量为 $60\ \mathrm{kg}$ 的人站在电梯中的体重计上，取 $g=10\ \mathrm{m/s^2}$。若电梯以 $2\ \mathrm{m/s^2}$ 的加速度向上启动，体重计示数对应的支持力是多少？

\textbf{物理模型：} 竖直方向牛顿第二定律，支持力 $N$ 就是视重。

\textbf{解题过程：} 取竖直向上为正方向，
\[
N-mg=ma.
\]
\[
N=m(g+a)=60\times(10+2)=720\ \mathrm{N}.
\]

\textbf{方法总结：} 电梯题先把体重计示数理解为支持力 $N$，不要当作重力。向上加速时 $N>mg$；若向下加速，同一公式中 $a$ 取负值。
\end{examplebox}

\section*{七、变式训练}

\begin{enumerate}
  \item 质量为 $m = 5\ \mathrm{kg}$ 的物体在水平面上，受水平拉力 $F = 25\ \mathrm{N}$，动摩擦因数 $\mu_k = 0.4$，$g = 10\ \mathrm{m/s^2}$。求加速度。
  
  \item 一个物体以 $v_0 = 10\ \mathrm{m/s}$ 的初速度沿倾角 $\theta = 30^\circ$ 的粗糙斜面向上滑行，动摩擦因数 $\mu_k = 0.2$，$g = 10\ \mathrm{m/s^2}$。求：
  \begin{center}
  \begin{tikzpicture}[scale=0.82,>=Stealth,line cap=round,line join=round]
  \coordinate (A) at (0,0);
  \coordinate (B) at (4.2,0);
  \coordinate (C) at (4.2,2.2);
  % 连续斜面先画在物块下方，避免右侧轮廓看成“小旗子”。
  \draw[thick] (A) -- (C) -- (B) -- cycle;
  \begin{scope}[shift={(2.15,1.13)},rotate=27.6]
    \draw[fill=gray!10,thick] (-0.45,0) rectangle (0.45,0.5);
  \end{scope}
  \draw[->,blue,thick] (1.72,1.1) -- (2.55,1.54) node[right] {$v_0$};
  \draw[->,red,thick] (2.18,1.29) -- (1.45,0.91) node[left] {$f_k$};
  \draw[->,red,thick] (2.15,1.05) -- (2.15,0.15) node[below] {$mg$};
  \draw[->,blue,thick] (2.15,1.05) -- (1.72,1.86) node[above left] {$N$};
  \draw (0.8,0) arc[start angle=0,end angle=27.6,radius=0.8];
  \node at (1.05,0.2) {$\theta$};
  \node[font=\small] at (3.12,2.12) {上滑过程};
\end{tikzpicture}

  \end{center}
  \begin{enumerate}
    \item 物体上滑的加速度大小和方向。
    \item 物体上滑的最大距离。
  \end{enumerate}
  
  \item 质量为 $50\ \mathrm{kg}$ 的人站在电梯中。若电梯以 $1.5\ \mathrm{m/s^2}$ 的加速度向下启动，取 $g=10\ \mathrm{m/s^2}$，求体重计示数对应的支持力。
  \begin{center}
  % TikZ 图：升降机中的超重、正常视重与失重
% Day 06：牛顿运动定律
\begin{tikzpicture}[scale=0.84,>=Latex,line cap=round,line join=round]
  % 向上加速：N 的箭头明显长于 mg。
  \begin{scope}[xshift=0cm]
    \node[font=\small] at (0,2.25) {向上加速};
    \draw[thick] (-1.0,-1.20) rectangle (1.0,1.20);
    \draw[thick] (-0.45,-0.75) rectangle (0.45,-0.10);
    \node at (0,-0.42) {$m$};
    \draw[->,very thick,blue] (0,-0.10) -- (0,1.02)
      node[right] {$N$};
    \draw[->,very thick,red] (0,-0.75) -- (0,-1.52)
      node[right] {$mg$};
    \draw[->,ultra thick] (-0.74,1.35) -- (-0.74,2.00)
      node[left] {$a$};
    \node[font=\small] at (0,-2.05) {$N>mg$};
  \end{scope}

  % 匀速或静止：两箭头等长。
  \begin{scope}[xshift=4cm]
    \node[font=\small] at (0,2.25) {匀速或静止};
    \draw[thick] (-1.0,-1.20) rectangle (1.0,1.20);
    \draw[thick] (-0.45,-0.75) rectangle (0.45,-0.10);
    \node at (0,-0.42) {$m$};
    \draw[->,very thick,blue] (0,-0.10) -- (0,0.78)
      node[right] {$N$};
    \draw[->,very thick,red] (0,-0.75) -- (0,-1.63)
      node[right] {$mg$};
    \node[font=\small] at (0,-2.05) {$N=mg$};
  \end{scope}

  % 向下加速：N 的箭头明显短于 mg。
  \begin{scope}[xshift=8cm]
    \node[font=\small] at (0,2.25) {向下加速};
    \draw[thick] (-1.0,-1.20) rectangle (1.0,1.20);
    \draw[thick] (-0.45,-0.75) rectangle (0.45,-0.10);
    \node at (0,-0.42) {$m$};
    \draw[->,very thick,blue] (0,-0.10) -- (0,0.52)
      node[right] {$N$};
    \draw[->,very thick,red] (0,-0.75) -- (0,-1.82)
      node[right] {$mg$};
    \draw[->,ultra thick] (-0.74,2.00) -- (-0.74,1.35)
      node[left] {$a$};
    \node[font=\small] at (0,-2.05) {$N<mg$};
  \end{scope}
\end{tikzpicture}

  \end{center}
\end{enumerate}

\section*{八、易错点}

\begin{mistakebox}
\begin{enumerate}
  \item \textbf{把 $F = ma$ 中的 $F$ 理解为某个力而不是合力：} 列式前先在选定方向上写出 $\sum F$，再令 $\sum F=ma$。物体在斜面上受重力、支持力、摩擦力，加速度 $a = g\sin\theta - \mu_k g\cos\theta$ 是由合力产生的，不是单由重力产生的。
  
  \item \textbf{混淆作用力与反作用力和平衡力：} 平衡力作用在同一物体上，作用力与反作用力作用在两个不同物体上。受力分析只画外界对研究对象的力，不画研究对象对外界的力。
  
  \item \textbf{正交分解时坐标方向不对：} 先观察加速度和约束方向。常把一个轴取在加速度方向或斜面方向上，另一个轴取垂直约束面方向，方程更清楚。
  
  \item \textbf{牛顿第二定律的瞬时性：} 遇到“刚断开”“刚撤去”时，先重画瞬间受力图；力一变，加速度立即改变。
  
  \item \textbf{连接体问题中内力不等于外力：} 先用整体法求共同加速度，再隔离某一物体求张力；不要把绳中张力直接等同于外部拉力。
  
  \item \textbf{把超重理解成重力变大：} 先把体重计示数写成支持力 $N$，再在竖直方向列 $N-mg=ma$。超重、失重改变的是支持力或悬绳拉力等视重，重力 $mg$ 在近地范围内近似不变。

  \item \textbf{有明显空气阻力时仍套自由落体：} 若阻力随速度增大，合力和加速度会随时间变化，不能直接用匀加速公式贯穿全过程。
\end{enumerate}
\end{mistakebox}

\section*{九、竞赛入门补充}

\begin{competitionbox}
\textbf{牛顿第二定律的瞬时性问题}

竞赛中常见“剪断绳子”、“撤去支持物”等瞬时变化问题。

\textbf{关键：} 弹簧弹力不能突变（形变需要时间），而绳的拉力可以突变。

\textbf{例题：} $A$、$B$ 两球用轻弹簧连接，$A$ 球用细绳悬挂在天花板上，如图。

\begin{center}
% TikZ 图：弹簧瞬时性问题
% Day 06：牛顿运动定律
% A球用细绳悬挂在顶，A和B用弹簧连接，剪断细绳瞬间
\begin{tikzpicture}[scale=0.9,>=Latex]
  % 天花板
  \draw[thick] (-2,0) -- (2,0);
  % 细绳悬挂A
  \draw[thick] (0,0) -- (0,-1.2);
  \draw[thick] (-0.5,-1.2) rectangle (0.5,-2.0);
  \node at (0,-1.6) {$A$};
  \node[left] at (0,-0.6) {细绳};
  % 弹簧连接A和B
  \draw[decorate,decoration={coil,aspect=0.3,segment length=3pt,amplitude=2pt}] (0,-2.0) -- (0,-3.3);
  \node[right] at (0,-2.65) {弹簧};
  % B球
  \draw[thick] (-0.5,-3.3) rectangle (0.5,-4.1);
  \node at (0,-3.7) {$B$};
  % 剪断示意
  \draw[red,thick] (0.6,-0.8) -- (-0.6,-1.0);
  \node[red,right] at (0.6,-0.9) {剪断};
  % 剪断后A受力
  \draw[->,thick,blue] (-1.2,-1.6) -- (-1.2,-2.4) node[left] {$G$};
  \draw[->,thick,red] (-1.2,-1.6) -- (-1.2,-0.8) node[left] {$T$（弹簧弹力不变）};
  % 剪断后B受力
  \draw[->,thick,blue] (1.2,-3.7) -- (1.2,-4.5) node[right] {$G$};
  \draw[->,thick,red] (1.2,-3.7) -- (1.2,-2.9) node[right] {$T$};
\end{tikzpicture}


\figcaption{剪断细绳瞬间：弹簧弹力 $T$ 保持不变，绳拉力 $F$ 立即消失}\label{fig:spring-instant}
\end{center}

剪断细绳的瞬间，$A$、$B$ 两球的加速度各是多少？

\textbf{解答：} 
\begin{itemize}
  \item 剪断前：$A$ 受重力 $m_A g$、弹簧拉力 $T$（向下）、绳拉力 $F$（向上），平衡 $F = m_A g + T$；$B$ 受重力 $m_B g$、弹簧拉力 $T$（向上），平衡 $T = m_B g$。
  \item 剪断瞬间：绳拉力 $F$ 立即消失，但弹簧形变未变，$T$ 保持不变（$T = m_B g$）。
  \item 对 $A$：$m_A g + T = m_A a_A$，$a_A = \dfrac{m_A + m_B}{m_A} g$（向下）。
  \item 对 $B$：$m_B g - T = m_B a_B$，$a_B = 0$（弹力与重力仍平衡）。
\end{itemize}

\textbf{结论：} 弹簧弹力在瞬间保持不变，绳子拉力可以突变。若 $m_A = m_B = m$，则剪断瞬间 $a_A = 2g$（向下），$a_B = 0$。

\medskip
\textbf{速度相关阻力与终端速度}

现实中的雨滴、降落伞和高速运动物体常受到空气阻力。阻力通常随速度增大而增大。为了建立基础模型，可以先假设阻力大小近似为
\[
f=kv,
\]
方向与速度相反。若物体竖直下落，取向下为正方向，则
\[
mg-kv=ma.
\]
刚释放时 $v=0$，阻力很小，加速度接近 $g$；速度增大后，阻力变大，合力变小，加速度也变小。当
\[
mg=kv_t
\]
时，加速度变为 0，物体达到稳定速度
\[
v_t=\frac{mg}{k},
\]
这叫\textbf{终端速度}。

\begin{center}
% TikZ 图：速度相关阻力与终端速度
% Day 06：牛顿第二定律扩展
\begin{tikzpicture}[>=Latex,scale=0.92,every node/.style={font=\scriptsize}]
  \begin{scope}
    \node[font=\small\bfseries] at (1.9,3.35) {受力变化};
    \draw[thick] (1.9,1.65) circle (0.28);
    \draw[blue,->,very thick] (1.72,1.35) -- (1.72,0.45) node[below] {$mg$};
    \draw[red,->,very thick] (2.08,1.95) -- (2.08,2.75);
    \node[red,right] at (2.08,2.55) {$f=kv$};
    \draw[->,thick] (2.75,1.95) -- (2.75,0.85) node[midway,right] {$v$};
  \end{scope}

  \begin{scope}[xshift=4.5cm]
    \node[font=\small\bfseries] at (2.1,3.15) {$v$-$t$ 图像};
    \draw[->,thick] (0,0) -- (4.25,0) node[right] {$t$};
    \draw[->,thick] (0,0) -- (0,2.85) node[above] {$v$};
    \draw[dashed] (0,2.15) -- (4.0,2.15);
    \node[left] at (0,2.15) {$v_t$};
    \draw[blue,very thick,domain=0:3.9,samples=80] plot(\x,{2.15*(1-exp(-0.9*\x))});
    \node[blue,align=center] at (2.55,1.20) {加速度逐渐减小};
    \node[align=center] at (2.55,2.55) {最终匀速\\$mg=kv_t$};
  \end{scope}
\end{tikzpicture}


\figcaption{速度相关阻力：速度越大阻力越大，最终可能达到终端速度}\label{fig:drag-terminal-velocity}
\end{center}

\textbf{模型边界：} 这里的 $f=kv$ 只是入门近似。真实空气阻力可能与 $v^2$ 更接近，也会受形状和空气密度影响。高中阶段的重点是理解：在“重力恒定、阻力随速度改变”的竖直下落模型中，合力会随速度改变，加速度不再恒定，不能机械套匀变速公式。

\pagebreak[4]
\medskip
\textbf{非惯性参考系与惯性力}

牛顿第二定律最直接适用于惯性参考系。若参考系本身正在加速，例如随一辆向右加速的车厢一起观察，车厢中的物体相对于车厢会表现为“向后偏”。这不是因为多了某个真实物体在向后拉它，而是因为观察者选了加速参考系。

\begin{center}
% TikZ 图：加速参考系中的惯性力
% Day 06：牛顿运动定律
\begin{tikzpicture}[scale=0.92,>=Latex]
  \tikzset{
    cart/.style={draw,thick,rounded corners=2pt,fill=gray!8},
    force/.style={->,very thick},
    note/.style={font=\small,align=center}
  }

  \begin{scope}[xshift=-3.3cm]
    \draw[cart] (-1.45,-0.75) rectangle (1.45,0.55);
    \draw[thick] (-1.05,-0.75) -- (-1.30,-1.10);
    \draw[thick] (1.05,-0.75) -- (1.30,-1.10);
    \draw[thick] (-1.55,-1.10) -- (1.55,-1.10);
    \draw[thick,fill=white] (-0.35,-0.55) rectangle (0.35,0.15);
    \node at (0,-0.20) {$m$};
    \draw[force,red] (0.35,-0.20) -- (1.10,-0.20) node[above] {$F$};
    \draw[force,blue] (-1.10,0.85) -- (-0.25,0.85) node[above] {$a$};
    \node[note] at (0,-1.55) {地面惯性系\\只画真实受力};
  \end{scope}

  \begin{scope}[xshift=3.3cm]
    \draw[cart] (-1.45,-0.75) rectangle (1.45,0.55);
    \draw[thick] (-1.05,-0.75) -- (-1.30,-1.10);
    \draw[thick] (1.05,-0.75) -- (1.30,-1.10);
    \draw[thick] (-1.55,-1.10) -- (1.55,-1.10);
    \draw[thick,fill=white] (-0.35,-0.55) rectangle (0.35,0.15);
    \node at (0,-0.20) {$m$};
    \draw[force,red] (0.35,-0.20) -- (1.10,-0.20) node[above] {$F$};
    \draw[force,orange] (-0.35,-0.20) -- (-1.10,-0.20) node[above] {$F_{\text{惯}}=ma_0$};
    \draw[force,blue] (-1.10,0.85) -- (-0.25,0.85) node[above] {$a_0$};
    \node[note] at (0,-1.55) {随车加速系\\补一个反向惯性力};
  \end{scope}

  \node[note] at (0,1.55) {同一问题可换参考系，但要先说明是否为惯性系};
\end{tikzpicture}


\figcaption{非惯性参考系：随车观察时可引入与参考系加速度相反的惯性力}\label{fig:noninertial-frame-force}
\end{center}

\textbf{使用规则：} 若参考系以加速度 $\vect{a}_0$ 相对惯性系运动，在这个非惯性系中可给每个质量为 $m$ 的物体补一个
\[
\vect{F}_{\text{惯}}=-m\vect{a}_0.
\]
这样可以把物体在非惯性系中的“平衡”问题写成
\[
\sum \vect{F}_{\text{真实}}+\vect{F}_{\text{惯}}=0.
\]

\textbf{模型边界：} 惯性力不是物体间真实相互作用力，没有反作用力对象。高中常规题优先在地面惯性系中分析；竞赛入门题中，加速车厢、转动圆盘和离心现象可以用惯性力作为辅助工具，但必须先说明参考系。
\end{competitionbox}

\section*{十、今日小结}

{\small
\begin{enumerate}[nosep,leftmargin=2em]
  \item 牛顿第一定律：力是改变运动状态的原因（不是维持运动的原因）。
  \item 牛顿第二定律 $F_{\text{合}} = ma$：受力分析 $\to$ 求合力 $\to$ 算加速度 $\to$ 运动学求解。
  \item 牛顿第三定律：作用力与反作用力等大反向、同生同灭，但作用在不同物体上。
  \item 正交分解法：先选方便的坐标轴，常把一轴取在加速度或运动方向上；若某方向无加速度，该方向合力为零。
  \item 连接体思路：整体法求加速度，隔离法求内力（Day 07 专题深入）。
  \item 超重、失重看视重变化，本质是竖直方向牛顿第二定律。
  \item 有速度相关阻力时，合力和加速度会随速度改变；终端速度来自合力为零。
  \item 牛顿第二定律默认在惯性参考系中使用；非惯性系可引入惯性力，但它不是一对真实相互作用力。
\end{enumerate}

\textbf{一分钟复盘：} 拿一道动力学题，只写流程不计算：研究对象是谁、坐标轴怎么取、每个方向的合力是什么、加速度指向哪里。流程清楚后再代数计算。
}
